\section{AI Agent Identity: Delegation en lugar de secretos compartidos}

La clase ya señalaba en 2025 la brecha de seguridad de los agent frameworks: los scripts y las service accounts llevan tiempo ejecutando tareas en nombre de personas o sistemas; el LLM agent solo hace que el comportamiento sea más dinámico, que haya más herramientas y que la intención expresada en lenguaje natural sea más difícil de acotar con precisión. El problema central no es «si el agent tiene nombre de usuario», sino quién lo autoriza, en nombre de quién actúa, a qué audience puede acceder, qué actions puede ejecutar, cuándo caduca y cómo se revoca.

\subsection{Principal, delegation y OAuth}

\term{service account} es una cuenta no humana que usan el software o una workload; \term{non-human identity} (NHI) designa en general a un principal de tipo service, workload, device o agent. \term{OAuth 2.0} es un delegated authorization framework que permite al client obtener un scoped access token sin tener que poseer la contraseña del usuario; \term{OpenID Connect} (OIDC) proporciona una capa de identidad sobre OAuth y expresa el resultado de la autenticación mediante un ID token.

\lecturefigure{12-agent-delegation.png}{El agent necesita un principal propio y una delegation acotada, en lugar de heredar todos los secretos del usuario o de la máquina.}{Subtítulos locales 24:10--27:50; OAuth token exchange; redibujo conceptual.}

\begin{knowledgebox}{Lectura de la figura: distinguir «quién es el agent» de «en nombre de quién actúa»}
El user/owner aporta la intent y el consent; el agent principal acredita la workload identity; el authorization server emite el token con base en la delegation, el scope, la audience y la expiry; después, el tool/resource server realiza el policy check y escribe el audit. \term{token exchange} permite al sistema intercambiar un token existente por uno nuevo dirigido a una audience específica y con permisos más reducidos, conservando la relación entre actor y subject, para no entregar directamente al agent el broad token del usuario.
\end{knowledgebox}

\teachervoice{McKinnon advierte a los estudiantes que no se dejen confundir por la idea de que «el agent es algo completamente nuevo»: los cron jobs, los scripts y las service accounts llevan mucho tiempo actuando en nombre de los sistemas. Lo verdaderamente nuevo es la ampliación de las capacidades y de la autonomía, mientras que muchos frameworks todavía dejan la security para el final e incluso permiten que las máquinas de desarrollo acumulen tokens con acceso a todos los sistemas.}

\subsection{El token es una capability con ciclo de vida}

\term{access token} es la credencial que expresa una autorización ante el resource server; un bearer token puede usarlo cualquiera que lo obtenga, por lo que es necesario limitar su scope, su audience y su lifetime, y adoptar en lo posible tokens sender-constrained/bound. Cuando la agent task termina, el owner deja la organización, el device presenta riesgo, el modelo se deshabilita o cambia la tool policy, debe ser posible hacer revoke de la session/token existente.

\lecturefigure{13-token-lifecycle.png}{El token es una capacidad de acceso temporal que requiere un ciclo de vida completo: issue, bind, use y revoke/retire.}{Conceptos de OAuth/OIDC/WebAuthn; redibujo conceptual.}

\begin{knowledgebox}{Lectura de la figura: el least privilege debe concretarse por tarea}
En la fase de issue se confirman el principal, la audience, el scope y la expiry; en la fase de bind se asocia el token a una key/device/session, lo que reduce su utilidad tras un token theft; en la fase de use, el resource server verifica la action y el trace ID; en la fase de revoke/retire se responde a risk, job completion u owner removal. Si un agent solo necesita leer un issue y crear un borrador, no debe recibir repo admin, billing ni un refresh token de larga duración.
\end{knowledgebox}

\begin{importantbox}{Campos mínimos de un delegation record}
Cada delegation a un agent debe registrar como mínimo: human/organizational owner, agent principal, subject, requested task, resource audience, approved scopes, policy/version, issue/expiry, token/session ID, tool actions, result y revocation reason. Solo así el incident response puede responder «qué persona autorizó a qué agent a hacer qué en qué sistema».
\end{importantbox}

\begin{warningbox}{La agent memory no debe convertirse en un credential warehouse}
El prompt, el memory store, la trace y el local workspace suelen diseñarse como contexto recuperable; si en ellos se mezclan API keys, refresh tokens o session cookies, las salidas del modelo, los complementos y los registros de depuración pueden ampliar la superficie de filtración. Los secrets deben guardarse en un vault dedicado, inyectarse por task mediante un short-lived credential broker y excluirse del contexto del modelo y de los registros ordinarios.
\end{warningbox}

\subsection{Resumen de la sección}

Lo esencial de la agent identity es la separación entre actor y subject, la delegation, el least privilege, el short-lived token y el audit. Compartir la contraseña del usuario o un token de larga duración convierte un solo compromise de un agent o de una máquina en un acceso persistente a varios sistemas.
